\documentclass[10pt,a4paper]{amsart}
\usepackage[margin=19mm]{geometry}
\usepackage{amsmath,amssymb,mathtools}
\usepackage[scr=rsfs,cal=euler]{mathalfa}
\usepackage{xfrac}
\usepackage[unicode,hidelinks,bookmarks=false]{hyperref}
\newcommand{\ie}{i.e.\ }
\newcommand{\cf}{cf.\ }
\newcommand{\resp}{respectively\ }
\newcommand{\Subsection}{\subsection}
\providecommand{\nobreakdash}{\nobreak\hbox{-}}
\setlength{\emergencystretch}{2em}
\multlinegap0pt
\usepackage{braket}
\usepackage{amssymb}
\usepackage{float}
\usepackage{tikz}
\newtheorem{cor}{Corollary}
\def\bra#1{\langle #1 |}
\def\ket#1{| #1 \rangle}
\title{Structure Theorem for Quantum Replacer Codes}
\author{Eric Chitambar, Sarah Hagen, David W. Kribs, Mike I. Nelson, Andrew Nemec}
\date{Source: arXiv:2505.06659v1 (CC BY 4.0)}
\begin{document}
\maketitle

\noindent\emph{Source and licence.} Structure Theorem for Quantum Replacer Codes. Version arXiv:2505.06659v1 (CC BY 4.0), distributed by arXiv under the Creative Commons Attribution 4.0 International licence, http://creativecommons.org/licenses/by/4.0/, as recorded in the arXiv licence metadata for this exact version and shown on the abstract page https://arxiv.org/abs/2505.06659v1. Full licence notice: \texttt{licenses/arxiv-2505.06659-CC-BY-4.0.txt}.\par\smallskip

\noindent\textit{Editorial scope.} Counted unit: the article's cor cleaninglemma (source0.tex lines 869--871). The article's complete original proof of it is delivered with it (source0.tex lines 873--899, one \texttt{proof} environment of 27 lines). No mathematics is written, repaired or re-derived: the statement and the proof are the article's own lines, in the article's own notation, and no retained line is substituted or re-flowed.  Retained uncounted as the source-local context the proof needs: lines 718--728.  Nothing was omitted from any retained span, so the package carries no omission record.  No retained line contains a citation, so the wrapper carries no bibliography and no bibliography entry.  No retained line contains a figure environment, an image-inclusion command or drawing code, so no image is delivered and the package declares no asset dependency.  Editorial formatting only, in addition: an AMS \texttt{amsart} preamble that restates the article's own declarations and macros used by the retained lines, namely the environments \texttt{cor} on separate counters, together with the standard packages those lines need.  The retained environments print in this document as Corollary 1 in order of appearance, whereas the article prints the same items with the numbering of its own full document.  The complete native source of the article as distributed in the arXiv e-print for this exact version is bundled unchanged as \texttt{sources/177898/source0.tex}.
\begin{cor}\label{correctablecor}
    Let $S$ be a subspace of $\mathcal H = \overline{E} E$. Then $S$ corrects erasure errors on $E$ if and only if for every operator $X = I_{\overline{E}} \otimes X_E$ on $\mathcal H$ that only acts non-trivially on $E$, there is a scalar $c(X)$ such that 
    \[
    \bra{\widetilde{\phi}} X \ket{\widetilde{\phi}} = c(X), 
    \]
    for all $\widetilde{\phi}\in S$. Further, when this condition is satisfied we have 
    \[
    c(X) =  \bra{\psi} \, (I_A \otimes X_E) \,   \ket{\psi} , 
    \]
    where $A$ is the ancilla and $\ket{\psi} = \ket{\psi}_{AE}$ is the state from condition~$(iv)$ of the structure theorem. 
\end{cor}

\begin{cor}\label{cleaninglemma}
Suppose $S$ is a subspace of $\mathcal H = \overline{E} E$ and let $P_S$ be the projection of $\mathcal H$ onto $S$. If $S$ is a correctable code for a replacer channel $\mathcal E_E$, then there are no unitary operators $L\in \{ P_S \}^\prime$ for which $L P_S \notin \mathbb{T} P_S$ and $\mathrm{supp}(L)\subseteq E$.
\end{cor}

\begin{proof}
Suppose we have a unitary operator $L\in \{ P_S \}^\prime$ such that $\mathrm{supp}(L)\subseteq E$, so that $L = I_{\overline{E}} \otimes L_E$ for some unitary $L_E$ on $E$. 

We can use the code state form given in condition~$(iv)$ of the structure theorem to find, for each basis state $ \ket{\widetilde{i}}$ of $S$, 
\begin{eqnarray*}
 L  \, \ket{\widetilde{i}} &=& \big( I_{\overline{E}} \otimes L_E \big) \big( U_{\overline{E}} \otimes I_E \big) \, ( \ket{i}_R \otimes \ket{\psi}_{A E}) \\ 
 &=&  \big( U_{\overline{E}} \otimes I_E \big) \big( I_{RA} \otimes L_E \big) \, ( \ket{i}_R \otimes \ket{\psi}_{A E}) \\ 
  &=&  \big( U_{\overline{E}} \otimes I_E \big) \, ( \ket{i}_R \otimes \big( ( I_{A} \otimes L_E   ) \ket{\psi}_{A E}\big) ) . 
\end{eqnarray*}
On the other hand, since $L$ commutes with $P_S$, there are scalars $\{a_j\}_j$ such that 
\[
L  \, \ket{\widetilde{i}} = L P_S \ket{\widetilde{i}} = P_S L   \ket{\widetilde{i}} = \sum_j a_j  \ket{\widetilde{j}} = \big( U_{\overline{E}} \otimes I_E \big) \, \big( \big( \sum_j a_j \ket{j}_R \big) \otimes \ket{\psi}_{A E} \big).  
\]

Hence we can multiply both descriptions of $L\ket{\widetilde{i}}$ by $U_{\overline{E}}^\dagger \otimes I_E$ and use the fact that $U_{\overline{E}}^\dagger U_{\overline{E}} = I_{RA}$ to obtain for all $ \ket{\widetilde{i}}$, 
\[
\ket{i}_R \otimes \big( ( I_{A} \otimes L_E   ) \ket{\psi}_{A E}\big) = 
\Big( \sum_j a_j \ket{j}_R \Big) \otimes \ket{\psi}_{A E} . 
\]
If we take the inner product of this vector with $\bra{j}_R\bra{\psi}_{AE}$ for any $j$, we get 
\[
\delta_{ij} \big( \,   \underbrace{\bra{\psi}_{AE} (I_{A} \otimes L_E   ) \ket{\psi}_{A E}}_{\lambda} \big) = a_j, 
\]
and note that the indicated scalar $\lambda$ is independent of $i$. (Observe this is the scalar of Corollary~\ref{correctablecor}.) 

Therefore,  we have $L \ket{\widetilde{i}} = \lambda \ket{\widetilde{i}}$ for all $i$, and so $LP_S = P_S L = \lambda P_S$ (also $\lambda\in\mathbb{T}$ as $L$ is unitary). The result follows.  
\end{proof}

\end{document}
